\documentclass[14pt]{extreport}
\usepackage{graphicx}
\usepackage{gost}
\usepackage[T1,T2A]{fontenc}
\usepackage[utf8]{inputenc}
\PassOptionsToPackage{russian,english}{babel}
\usepackage[english,russian]{babel}
% \usepackage{tempora}
\usepackage{hyperref}
\linespread{1.3}
\setlength{\parindent}{1.25cm}
\usepackage{hyperref}

\usepackage{multirow}
\usepackage{makecell}
\usepackage{longtable}
\usepackage{array}
\newcolumntype{P}[1]{>{\raggedright\arraybackslash}p{#1}}

\usepackage{fancyhdr}
\pagestyle{fancy}
\fancyhf{} 
\fancyhead{} 
\fancyfoot[C]{\thepage} 

\usepackage{amsmath, amssymb}
\usepackage{tcolorbox}
\usepackage{enumitem}
\usepackage{geometry}
\usepackage{amsthm}

\begin{document}
\thispagestyle{empty}
	\title{Математический анализ. Экзамен.}
	\maketitle

	\newpage
	\tableofcontents

    \newpage
    \chapter{Понятие первообразной. Свойства первообразной. Неопределенный интеграл. Свойства неопределенного интеграла.}

    \begin{definition}
        Функция $F(x)$ называется \textbf{первообразной} функции $f(x)$, если $F'(x) = f(x)$
    \end{definition}

    \begin{theorem}[Общий вид первообразной]
        Пусть $F_1(x) $ и $F_2(x)$ -- первообразные функции $f(x)$. Тогда $F_1(x) - F_2(x) = C$
    \end{theorem}
    \begin{proof}
        Пусть $H(x) = F_2(x) - F_1(x)$. $H'(x) = f(x) - f(x) = 0$. Докажем, что если $H'(x) = 0$, то $H(X) = const$. Пусть $H(X)$ определена на $[a, b]$, $(a, b) \subset D(f)$. $H(x)$ непрерывна на $[a, b] \Rightarrow H(x)$ непрерывна на $[a_1, b_1] \subset [a, b]$. Тогда применим теорему Лагранжа к $H(x): \exists c \in (a_1, b_1) : H(c) = \frac{H(b_1)-H(a_1)}{b_1-a_1}$. $\forall c \: H'(c)=0 \Rightarrow H(a_1)=H(b_1)\Rightarrow H(x )=const$
    \end{proof}
    \begin{center}
        \renewcommand{\arraystretch}{1.5}
        \begin{tabular}{|c|c|||c|c|}
        \hline
        \textbf{Функция} & \textbf{Интеграл} & \textbf{Функция} & \textbf{Интеграл} \\
        \hline
        $0$ & $C$ & $\cos x$ & $\sin x + C$ \\
        \hline
        $1$ & $x + C$ & $\sin x$ & $-\cos x + C$ \\
        \hline
        $x^n \ (n \neq -1)$ & $\frac{x^{n+1}}{n+1} + C$ & $\frac{1}{\cos^2 x}$ & $\operatorname{tg} x + C$ \\
        \hline
        $\frac{1}{x}$ & $\ln|x| + C$ & $\frac{1}{\sin^2 x}$ & $-\operatorname{ctg} x + C$ \\
        \hline
        $a^x \ (a>0, a\neq 1)$ & $\frac{a^x}{\ln a} + C$ & $\frac{1}{\sqrt{1-x^2}}$ & $\arcsin x + C$ \\
        \hline
        $e^x$ & $e^x + C$ & $\frac{1}{1+x^2}$ & $\operatorname{arctg} x + C$ \\
        \hline
        \end{tabular}
    \end{center}

    \begin{definition}
        Множество всех первообразных функций $f(x)$ нахывается \textbf{неопределенным интегралом} от функции $f$
        $$
        F(x) = \int f(x)dx
        $$
    \end{definition}
    \newpage
    \begin{theorem}
        Неопределенный интеграл обладает следующими основными свойствами:
        \begin{enumerate}
            \item $\int f(x)dx + \int g(x)dx = \int (f(x) + g(x)) dx$
            \item $\int f(g(x)) g'(x) dx = \int f(g(x)) dg(x) = \int f(t) dt = F(g(x))$
            \item $\int f(ax+b) dx = \frac{1}{a} F(ax+b)+C$
            \item $\frac{1}{a}\int f(ax+b) \cdot a dx=\frac{1}{a} \int f(ax+b)(ax+b)'dx=\frac{1}{a}\int f(t) dt = \frac{1}{a}F(t)+C$
        \end{enumerate}
    \end{theorem}

    \chapter{Таблица неопределённых интегралов. Непосредственное интегрирование.}
    \section{Таблица основных неопределённых интегралов}

    \begin{center}
        \renewcommand{\arraystretch}{1.5}
        \begin{tabular}{|c|c|||c|c|}
        \hline
        \textbf{Функция} & \textbf{Интеграл} & \textbf{Функция} & \textbf{Интеграл} \\
        \hline
        $0$ & $C$ & $\cos x$ & $\sin x + C$ \\
        \hline
        $1$ & $x + C$ & $\sin x$ & $-\cos x + C$ \\
        \hline
        $x^n \ (n \neq -1)$ & $\frac{x^{n+1}}{n+1} + C$ & $\frac{1}{\cos^2 x}$ & $\operatorname{tg} x + C$ \\
        \hline
        $\frac{1}{x}$ & $\ln|x| + C$ & $\frac{1}{\sin^2 x}$ & $-\operatorname{ctg} x + C$ \\
        \hline
        $a^x \ (a>0, a\neq 1)$ & $\frac{a^x}{\ln a} + C$ & $\frac{1}{\sqrt{1-x^2}}$ & $\arcsin x + C$ \\
        \hline
        $e^x$ & $e^x + C$ & $\frac{1}{1+x^2}$ & $\operatorname{arctg} x + C$ \\
        \hline
        \end{tabular}
    \end{center}

    \section{Непосредственное интегрирование}

    \begin{definition}
        \textbf{Непосредственным интегрированием} называется вычисление неопределённого интеграла путём применения таблицы интегралов и свойств неопределённого интеграла без использования специальных методов (замены переменной, интегрирования по частям и др.).
    \end{definition}

    Основой непосредственного интегрирования являются следующие свойства:

    \begin{theorem}[Линейность неопределённого интеграла]
    Пусть функции $f(x)$ и $g(x)$ имеют первообразные, а числа
    $\alpha,\beta\in\mathbb R$. Тогда
    \[
    \int \bigl(\alpha f(x)+\beta g(x)\bigr)\,dx
    =
    \alpha\int f(x)\,dx+
    \beta\int g(x)\,dx.
    \]
    \end{theorem}

    \begin{proof}
    Пусть
    \[
    F'(x)=f(x), \qquad G'(x)=g(x).
    \]
    Тогда
    \[
    (\alpha F(x)+\beta G(x))'
    =
    \alpha F'(x)+\beta G'(x)
    =
    \alpha f(x)+\beta g(x).
    \]
    Следовательно,
    \[
    \alpha F(x)+\beta G(x)
    \]
    является первообразной функции
    \[
    \alpha f(x)+\beta g(x),
    \]
    откуда
    \[
    \int (\alpha f(x)+\beta g(x))\,dx
    =
    \alpha\int f(x)\,dx
    +
    \beta\int g(x)\,dx.
    \]
    \end{proof}

    \chapter{Дифференциал, инвариантность формы дифференциала. Интегрирование методом замены переменной в неопределенном интеграле.}

    Пусть функция $y=f(x)$ дифференцируема в точке $x$.

    \begin{definition}
    Дифференциалом функции $y=f(x)$ называется главная линейная часть её приращения:
    \[
    dy=f'(x)\,dx,
    \]
    где $dx=\Delta x$ называется дифференциалом независимой переменной.
    \end{definition}

    \section{Свойства дифференциала}

Если функции $u(x)$ и $v(x)$ дифференцируемы, то:

\begin{enumerate}
\item
\[
d(C)=0.
\]

\item
\[
d(Cu)=C\,du.
\]

\item
\[
d(u\pm v)=du\pm dv.
\]

\item
\[
d(uv)=u\,dv+v\,du.
\]

\item
\[
d\left(\frac{u}{v}\right)
=
\frac{v\,du-u\,dv}{v^2},
\qquad v\neq0.
\]
\end{enumerate}

Все свойства непосредственно следуют из соответствующих правил
дифференцирования.

\section{Инвариантность формы дифференциала}

\begin{theorem}
Форма записи первого дифференциала не изменяется при замене независимой
переменной.
\end{theorem}

\begin{proof}

Пусть

\[
y=f(u),
\qquad
u=\varphi(x),
\]

причём обе функции дифференцируемы.

Тогда по правилу дифференцирования сложной функции

\[
\frac{dy}{dx}
=
\frac{dy}{du}\cdot\frac{du}{dx}.
\]

Умножая на $dx$, получаем

\[
dy
=
\frac{dy}{du}\frac{du}{dx}\,dx.
\]

Так как

\[
du=\frac{du}{dx}\,dx,
\]

то

\[
dy
=
\frac{dy}{du}\,du.
\]

Следовательно, независимо от того, является ли аргументом функции
переменная $x$ или промежуточная переменная $u$, дифференциал имеет одну и
ту же форму:

\[
dy=f'(u)\,du.
\]

Теорема доказана.
\end{proof}

\section{Замена переменной в неопределённом интеграле}

\begin{theorem}[Замена переменной]
Если $x=\varphi(t)$
дифференцируема и после подстановки интеграл принимает вид
\[
\int f(\varphi(t))\varphi'(t)\,dt,
\]
то
\[
\int f(x)\,dx
=
\int f(\varphi(t))\varphi'(t)\,dt.
\]
\end{theorem}

\begin{proof}

Пусть $F'(x)=f(x).$ Тогда

\[
\frac{d}{dt}F(\varphi(t))
=
F'(\varphi(t))\varphi'(t)
=
f(\varphi(t))\varphi'(t).
\]

Следовательно,

\[
\int f(\varphi(t))\varphi'(t)\,dt
=
F(\varphi(t))+C.
\]

Возвращаясь к переменной $x$:

\[
F(\varphi(t))
=
F(x).
\]

Поэтому

\[
\int f(x)\,dx
=
F(x)+C.
\]

Теорема доказана.
\end{proof}
\end{document}